% from ega4-17-fr.tex
\oldpage[IV-4]{70}
en ce point toutes les propriétés définies dans (6.8.1). On a donc les propriétés suivantes, que nous rappelons pour la commodité des références:

\begin{proposition}[17.5.7]
\label{IV.17.5.7}
Soit $f:X\to Y$ un morphisme localement de présentation finie, lisse en un point $x\in X$; posons $y=f(x)$. Alors, pour que l'anneau $\sh{O}_{X,x}$ soit réduit (resp. intégralement clos, resp. géométriquement unibranche), il faut et il suffit que $\sh{O}_{Y,y}$ le soit.
\end{proposition}

Cela a en effet été prouvé dans (11.3.13) et (11.3.14) complété par Err$_{\mathrm{IV}}$, 30.

\begin{proposition}[17.5.8]
\label{IV.17.5.8}
Soient $Y$ un préschéma localement noethérien, $f:X\to Y$ un morphisme localement de type fini, lisse en un point $x\in X$; posons $y=f(x)$. Alors:
\begin{enumerate}
  \item[(i)] On a $\dim(\sh{O}_{X,x})=\dim(\sh{O}_{Y,y})+\dim(\sh{O}_{X,x}\otimes_{\sh{O}_{Y,y}}k(y))$.
  \item[(ii)] On a $\operatorname{coprof}(\sh{O}_{X,x})=\operatorname{coprof}(\sh{O}_{Y,y})$.
  \item[(iii)] Pour que l'anneau $\sh{O}_{X,x}$ possède la propriété $(S_n)$ (5.7.2) (resp. $(R_n)$ (5.8.2)), il faut et il suffit que l'anneau $\sh{O}_{Y,y}$ la possède. En particulier, pour que $\sh{O}_{X,x}$ soit régulier, il faut et il suffit que $\sh{O}_{Y,y}$ le soit.
\end{enumerate}
\end{proposition}

Ce sont des cas particuliers de (6.1.2), (6.3.2), (6.4.1) et (6.5.3).

\subsection{Caractérisations des morphismes étales}
\EGAReaderTarget{ega4.c04.subsec.17-6}{subsection:IV.17.6}

\begin{theorem}[17.6.1]
\label{IV.17.6.1}
Soient $f:X\to Y$ un morphisme localement de présentation finie, $x$ un point de $X$, $y=f(x)$. Les conditions suivantes sont équivalentes:
\begin{enumerate}
  \item[a)] $f$ est étale au point $x$.
  \item[a')] $f$ est lisse au point $x$ et non ramifié au point $x$.
  \item[b)] $f$ est lisse au point $x$ et quasi-fini au point $x$ (Err$_{\mathrm{III}}$, 20).
  \item[c)] $f$ est plat au point $x$ et non ramifié au point $x$.
  \item[c')] $f$ est plat au point $x$ et l'anneau $\sh{O}_{X,x}/\mathfrak m_y\sh{O}_{X,x}$ est un corps, extension finie séparable de $k(y)$.
  \item[d)] L'anneau $\sh{O}_{X,x}$ est une $\sh{O}_{Y,y}$-algèbre formellement étale pour les topologies discrètes.
\end{enumerate}
\end{theorem}

L'équivalence de a) et a') résulte aussitôt des définitions; celle de a) et d) résulte de l'équivalence de a) et e) dans (17.4.1) et de l'équivalence de a) et d) dans (17.5.1). L'équivalence de c) et c') résulte de l'équivalence de a) et d') dans (17.4.1). Le fait que a') implique c') découle de (17.5.1); inversement, si c') est vérifiée, $f$ est régulier (donc lisse par (17.5.1)) au point $x$, car si $K$ est une extension finie séparable d'un corps $k$, alors, pour toute extension $k'$ de $k$, $\Spec(K\otimes_k k')$ est régulier, étant somme d'un nombre fini de spectres de corps. Le fait que a') implique b) résulte de (17.4.1, d')) et de (17.5.1, b)). Il reste donc à voir que b) entraîne c), et comme on sait déjà que $f$ est plat au point $x$, par (17.5.1), il suffit de montrer que $f^{-1}(y)$ est un $k(y)$-préschéma non ramifié sur $k(y)$; autrement dit, on est ramené à prouver que b) entraîne c) lorsque $Y=\Spec(k)$ est le spectre d'un corps $k$. Comme la question est locale sur $X$, on peut se borner au cas où $X=\Spec(A)$, où $A$ est une $k$-algèbre finie et locale ($0_{\mathrm I}$, 7.4.1). En vertu de l'hypothèse b), $A$ est une $k$-algèbre formellement lisse pour les topologies discrètes, qui coïncident ici avec les topologies préadiques; donc (0, 19.6.5), $A$ est un anneau local régulier, donc

